% Abstract for the ESIPS honours thesis.
% Intended to be pulled into main.tex via % Abstract for the ESIPS honours thesis.
% Intended to be pulled into main.tex via % Abstract for the ESIPS honours thesis.
% Intended to be pulled into main.tex via % Abstract for the ESIPS honours thesis.
% Intended to be pulled into main.tex via \input{tex_files/0. Abstract/abstract.tex}.
% Assumes an IEEEtran document; wrapped in the abstract environment so it renders
% in the single-column header before \maketitle's body. Contains no preamble.
\begin{abstract}
Deployed literature-search engines force a choice between structured retrieval,
which answers a query deterministically and auditably, and semantic retrieval,
which recovers relevance that exact terms miss. This thesis reconciles the two
paradigms in an industry setting, designing, building, and evaluating a
production-ready scholarly search system for Latent Knowledge, whose users are
biomedical researchers. A typed knowledge graph of approximately 150 million
biomedical works, ingested from the 2026-02-03 OpenAlex snapshot, anchors a
three-stage retrieval cascade of structured graph retrieval, bi-encoder dense
retrieval, and cross-encoder re-ranking. A language model is confined to the two
artefacts the structured stage requires, namely the per-work keyword index built
offline and the intermediate representation translated from each natural-language
query, and a deterministic compiler turns that representation into Cypher, so
that query construction never leaves the compiler. Component benchmarks select
every deployed model. The translation model reaches a joint accuracy of 0.937 on
a hand-built stratified gold set of 116 queries, the deployed keyword extractor
produces the strongest benchmarked index of any candidate across three
scientific keyphrase datasets, the leading bi-encoders and cross-encoders form
statistical ties at averages around 82 to 83 nDCG on SciRepEval ad-hoc search,
and quantising the stored embeddings to int8 is measurably free, with 99.8
percent of ranking quality retained at half the stored footprint. The comparison
of the assembled cascade against a retrieval-augmented LLM and Google Scholar is
specified in the research design and reported in the final submission.
\end{abstract}.
% Assumes an IEEEtran document; wrapped in the abstract environment so it renders
% in the single-column header before \maketitle's body. Contains no preamble.
\begin{abstract}
Deployed literature-search engines force a choice between structured retrieval,
which answers a query deterministically and auditably, and semantic retrieval,
which recovers relevance that exact terms miss. This thesis reconciles the two
paradigms in an industry setting, designing, building, and evaluating a
production-ready scholarly search system for Latent Knowledge, whose users are
biomedical researchers. A typed knowledge graph of approximately 150 million
biomedical works, ingested from the 2026-02-03 OpenAlex snapshot, anchors a
three-stage retrieval cascade of structured graph retrieval, bi-encoder dense
retrieval, and cross-encoder re-ranking. A language model is confined to the two
artefacts the structured stage requires, namely the per-work keyword index built
offline and the intermediate representation translated from each natural-language
query, and a deterministic compiler turns that representation into Cypher, so
that query construction never leaves the compiler. Component benchmarks select
every deployed model. The translation model reaches a joint accuracy of 0.937 on
a hand-built stratified gold set of 116 queries, the deployed keyword extractor
produces the strongest benchmarked index of any candidate across three
scientific keyphrase datasets, the leading bi-encoders and cross-encoders form
statistical ties at averages around 82 to 83 nDCG on SciRepEval ad-hoc search,
and quantising the stored embeddings to int8 is measurably free, with 99.8
percent of ranking quality retained at half the stored footprint. The comparison
of the assembled cascade against a retrieval-augmented LLM and Google Scholar is
specified in the research design and reported in the final submission.
\end{abstract}.
% Assumes an IEEEtran document; wrapped in the abstract environment so it renders
% in the single-column header before \maketitle's body. Contains no preamble.
\begin{abstract}
Deployed literature-search engines force a choice between structured retrieval,
which answers a query deterministically and auditably, and semantic retrieval,
which recovers relevance that exact terms miss. This thesis reconciles the two
paradigms in an industry setting, designing, building, and evaluating a
production-ready scholarly search system for Latent Knowledge, whose users are
biomedical researchers. A typed knowledge graph of approximately 150 million
biomedical works, ingested from the 2026-02-03 OpenAlex snapshot, anchors a
three-stage retrieval cascade of structured graph retrieval, bi-encoder dense
retrieval, and cross-encoder re-ranking. A language model is confined to the two
artefacts the structured stage requires, namely the per-work keyword index built
offline and the intermediate representation translated from each natural-language
query, and a deterministic compiler turns that representation into Cypher, so
that query construction never leaves the compiler. Component benchmarks select
every deployed model. The translation model reaches a joint accuracy of 0.937 on
a hand-built stratified gold set of 116 queries, the deployed keyword extractor
produces the strongest benchmarked index of any candidate across three
scientific keyphrase datasets, the leading bi-encoders and cross-encoders form
statistical ties at averages around 82 to 83 nDCG on SciRepEval ad-hoc search,
and quantising the stored embeddings to int8 is measurably free, with 99.8
percent of ranking quality retained at half the stored footprint. The comparison
of the assembled cascade against a retrieval-augmented LLM and Google Scholar is
specified in the research design and reported in the final submission.
\end{abstract}.
% Assumes an IEEEtran document; wrapped in the abstract environment so it renders
% in the single-column header before \maketitle's body. Contains no preamble.
\begin{abstract}
Deployed literature-search engines force a choice between structured retrieval,
which answers a query deterministically and auditably, and semantic retrieval,
which recovers relevance that exact terms miss. This thesis reconciles the two
paradigms in an industry setting, designing, building, and evaluating a
production-ready scholarly search system for Latent Knowledge, whose users are
biomedical researchers. A typed knowledge graph of approximately 150 million
biomedical works, ingested from the 2026-02-03 OpenAlex snapshot, anchors a
three-stage retrieval cascade of structured graph retrieval, bi-encoder dense
retrieval, and cross-encoder re-ranking. A language model is confined to the two
artefacts the structured stage requires, namely the per-work keyword index built
offline and the intermediate representation translated from each natural-language
query, and a deterministic compiler turns that representation into Cypher, so
that query construction never leaves the compiler. Component benchmarks select
every deployed model. The translation model reaches a joint accuracy of 0.937 on
a hand-built stratified gold set of 116 queries, the deployed keyword extractor
produces the strongest benchmarked index of any candidate across three
scientific keyphrase datasets, the leading bi-encoders and cross-encoders form
statistical ties at averages around 82 to 83 nDCG on SciRepEval ad-hoc search,
and quantising the stored embeddings to int8 is measurably free, with 99.8
percent of ranking quality retained at half the stored footprint. The comparison
of the assembled cascade against a retrieval-augmented LLM and Google Scholar is
specified in the research design and reported in the final submission.
\end{abstract}